% !TeX spellcheck = en_US
\addscenariosection{2}{ }{Stacking Units}{\images/population_token.png}

\begin{multicols*}{2}

\subsection*{\MakeUppercase{Stacking Rules}}

Variant rules for stacking your units above ``Pack''.

When Stacking a ``Few'' Unit, add a marker on that Unit. The Stacked Unit has the same stats as the ``Few'' Unit. When the Stacked Unit's Health goes to 0, remove the Stack Token and add any excess \svg{damage} to the ``Few'' side.

``Few'' Units: maximum 1 Stack Token total.

When Stacking a ``Pack'' Unit, add a Stack Token for either +1 \svg{attack}, +1 \svg{defense}, +1 \svg{health_points}, or +2 \svg{initiative} on that Unit. The Stacked ``Pack'' Unit has the stats from the Unit Card plus bonuses from Stack Token(s). When the Stacked Unit's Health goes to 0, remove the Stack Token and add any excess \svg{damage} to the Unit.

``Pack'' Units: maximum 1 Stack Token of each type. If Stack Tokens run out, use something else to mark which stat is boosted.

On every Resource Round, add a Cube to all your built Dwellings. You may use any Cubes on a Dwelling to Stack that Dwelling's Units once per Cube.

For Stacking Costs, treat Units as Tiers: \svg{bronze} Units are T1--T3, \svg{silver} Units are T4--T5, and \svg{golden} Units are T6--T7. X is the Unit's Tier level.

\subsection*{\MakeUppercase{Stacking ``Few'' Units}}
\begin{itemize}
  \item Recruitment cost + X × \svg{gold}
\end{itemize}

\subsection*{\MakeUppercase{Stacking ``Pack'' Units}}
\begin{itemize}
  \item Reinforcement cost + X \svg{gold} × the number of Stack Tokens on that Unit after Stacking.
\end{itemize}

\columnbreak

\noindent
\includegraphics[
  width=\linewidth,
  height=0.85\textheight,
  keepaspectratio
]{\art/elf.jpg}

\end{multicols*}
